% Generated draft: requires mathematical and source review.
\documentclass[11pt]{article}
\usepackage[margin=25mm]{geometry}
\usepackage{fontspec}
\setmainfont{Latin Modern Roman}
\usepackage{amsmath,amssymb,mathtools,unicode-math}
\setmathfont{Latin Modern Math}
\usepackage{enumitem,needspace,xurl,hyperref}
\setlength{\parindent}{0pt}
\setlength{\parskip}{5pt}
\setlength{\emergencystretch}{3em}
\newcommand{\N}{\mathbb{N}}
\newcommand{\Z}{\mathbb{Z}}
\newcommand{\Q}{\mathbb{Q}}
\newcommand{\R}{\mathbb{R}}
\newcommand{\C}{\mathbb{C}}
\newcommand{\F}{\mathbb{F}}
\newcommand{\A}{\mathbb{A}}
\newcommand{\PP}{\mathbb{P}}
\newcommand{\E}{\mathbb{E}}
\newcommand{\eps}{\varepsilon}
\newcommand{\dd}{\,\mathrm{d}}
\begin{document}
\section*{UnsolvedMath ID: 701546}
\textbf{Source problem code:} catalogue entry
\subsection{Definitions and mathematical statement}
A profinite group is a topological group isomorphic to an inverse limit of finite discrete groups. A group is \emph{periodic} if every element has finite order; no uniform bound on these orders is assumed. For a profinite group \(G\), let
\[
Z(G)=\{z\in G:zg=gz\text{ for every }g\in G\}
\]
be its center, which is closed. Thus the hypothesis that \(G/Z(G)\) is periodic means
\[
\forall g\in G\;\exists n\in\Z_{\geq 1}\quad g^n\in Z(G).
\]
Define the abstract derived subgroup by
\[
D(G)=\langle x^{-1}y^{-1}xy:x,y\in G\rangle,
\]
and let \(\overline{D(G)}\) denote its closure in \(G\).

The supplied question, proposed by P. Shumyatsky, asks whether \([G,G]\) is necessarily periodic whenever \(G\) is profinite and \(G/Z(G)\) is periodic. Under the abstract derived subgroup interpretation \([G,G]=D(G)\), the requested assertion is
\[
\forall\,G\text{ profinite},\qquad
\left(\forall g\in G\;\exists n\in\Z_{\geq1}\;g^n\in Z(G)\right)
\Longrightarrow
\left(\forall d\in D(G)\;\exists m\in\Z_{\geq1}\;d^m=1\right).
\]
The supplied record does not specify whether \([G,G]\) means the abstract derived subgroup or the closed derived subgroup. Under the latter convention, the conclusion instead quantifies over every \(d\in\overline{D(G)}\). This convention requires clarification from the source; the two formulations are not identified here.
\subsection{Short English statement}
For a profinite group whose quotient by its center consists entirely of elements of finite order, must every element of its derived subgroup also have finite order? The supplied notation leaves unresolved whether the derived subgroup includes its topological closure.
\subsection{Sources}
\begin{itemize}
\item Local UnsolvedMath registry: supplied complete source record, permanent ID \(701546\), exact source problem code \texttt{catalogue entry}, title \emph{Kourovka Notebook Problem 18.118}. The record labels the problem \emph{open}; this is supplied registry status, not an independently verified assessment.
\item Supplied source attribution: \emph{The Kourovka Notebook (Unsolved Problems in Group Theory)}, original problem 18.118, proposed by P. Shumyatsky in 2014. The background identifies the 18th issue (2014) and the 21st issue edition (2026), arXiv identifier \texttt{1401.0300v48}. The supplied collection description names E. I. Khukhro and V. D. Mazurov as editors.
\end{itemize}
\end{document}
